% Generated draft: requires mathematical and source review.
\documentclass[11pt]{article}
\usepackage[margin=25mm]{geometry}
\usepackage{fontspec}
\setmainfont{Latin Modern Roman}
\usepackage{amsmath,amssymb,mathtools,unicode-math}
\setmathfont{Latin Modern Math}
\usepackage{enumitem,needspace,xurl,hyperref}
\setlength{\parindent}{0pt}
\setlength{\parskip}{5pt}
\setlength{\emergencystretch}{3em}
\newcommand{\N}{\mathbb{N}}
\newcommand{\Z}{\mathbb{Z}}
\newcommand{\Q}{\mathbb{Q}}
\newcommand{\R}{\mathbb{R}}
\newcommand{\C}{\mathbb{C}}
\newcommand{\F}{\mathbb{F}}
\newcommand{\A}{\mathbb{A}}
\newcommand{\PP}{\mathbb{P}}
\newcommand{\E}{\mathbb{E}}
\newcommand{\eps}{\varepsilon}
\newcommand{\dd}{\,\mathrm{d}}
\begin{document}
\section*{UnsolvedMath ID: 701202}
\textbf{Source problem code:} catalogue entry
\subsection{Definitions and mathematical statement}
Under the customary interpretation of the terminology, let \(F_n=F(X)\) be the free group on a specified finite basis \(X=\{x_1,\ldots,x_n\}\), with \(n\geq 1\). An element \(p\in F_n\) is primitive if it belongs to some free basis of \(F_n\). Define the primitive length of \(g\in F_n\) by
\[
\ell_{\mathrm{prim}}(g)=\min\left\{k\in\Z_{\geq 0}: g=p_1\cdots p_k,\quad p_i\text{ is primitive for every }1\leq i\leq k\right\},
\]
where the empty product is the identity, so that \(\ell_{\mathrm{prim}}(1)=0\).

The question is whether there exists an algorithm which, given an element \(g\in F_n\) as a finite reduced word in \(X\cup X^{-1}\), terminates and returns the exact integer \(\ell_{\mathrm{prim}}(g)\).

The source refers to Problem 14.84 for the definition of a primitive element and to Problem 16.9 for the analogous definition of palindromic length; neither referenced text is supplied. Consequently, the definitions above, including the empty-product convention, require comparison with those references. The source also does not specify the intended range of ranks or whether the algorithm is sought separately for each fixed rank or uniformly with the rank included in the input. The assumption \(n\geq 1\) describes the customary setting of this draft and does not resolve that missing rank qualification.
\subsection{Short English statement}
Can one compute the minimum number of primitive elements whose product is a given element of a finite-rank free group? Here primitive means belonging to a free basis. The source's rank quantification and referenced definitions remain to be checked.
\subsection{Sources}
\begin{itemize}
\item Supplied local UnsolvedMath registry record: permanent ID \(701202\), exact source code \texttt{catalogue entry}, title \emph{Kourovka Notebook Problem 16.10}.
\item \emph{The Kourovka Notebook (Unsolved Problems in Group Theory)}, original Problem 16.10, 16th issue (2006); the supplied record identifies the 21st issue edition (2026), \texttt{arXiv:1401.0300v48}. The collection description names E. I. Khukhro and V. D. Mazurov as editors.
\item Proposed by V. G. Bardakov, V. A. Tolstykh, and V. E. Shpilrain; proposed year supplied as 2006.
\item The statement explicitly refers to Problems 14.84 and 16.9 for definitions. Their contents were not supplied.
\item The local record labels the problem \emph{open} and gives an update timestamp of \texttt{2026-10-16T14:00:00Z}. This is supplied registry information, not an independently verified mathematical or literature assessment.
\end{itemize}
\end{document}
